\subsection{Radon-Nikodym Theorem}

\begin{theorem} \label{thm:radon-nikodym-sigma-finite} \hlt{Radon--Nikodym Theorem}\\
Let $\lambda,\mu$ be $\sigma$-finite measures on $(X,\mathcal{M})$ with $\lambda\ll\mu$. Then there is a measurable $f:X\to[0,+\infty]$, unique $\mu$-a.e. and finite $\mu$-a.e., such that $\lambda(E)=\int_Ef\,d\mu$ for all $E\in\mathcal{M}$. Write $f=d\lambda/d\mu$.
\begin{proof}
Assume first $\lambda(X),\mu(X)<+\infty$.\\
For $c>0$, let $(P_k,N_k)$ be Hahn for $\lambda-kc\mu$, put $A_k=N_k\setminus\bigcup_{j<k}A_j$, and $B=X\setminus\bigcup_{k\geq1}A_k$.\\
Then $A_k\subseteq N_k$ and, for $k\geq2$, $A_k\subseteq P_{k-1}$; hence $E\subseteq A_k$ implies $(k-1)c\,\mu(E)\leq\lambda(E)\leq kc\,\mu(E)$.\\
Also $B\subseteq P_k$ for all $k$, so $\mu(B)=\lambda(B)=0$.\\
Define $f_c\defeq\sum_{k\geq1}(k-1)c\,\mathbf{1}_{A_k}$. Then $\int_E f_c\,d\mu\leq\lambda(E)\leq\int_E f_c\,d\mu+c\mu(X)$ for all $E\in\mathcal{M}$.\\
With $c_n=2^{-n}$ and $f_n=f_{c_n}$, $\norm{f_n-f_m}_1\leq(2^{-n}+2^{-m})\mu(X)\leq2^{-n+1}\mu(X)$ for $m\geq n$, so $f_n\to f$ in $L^1$ for some $f\geq0$. Letting $n\to\infty$ gives $\lambda(E)=\int_Ef\,d\mu$.\\
For $\sigma$-finite $\lambda,\mu$, choose disjoint $Y_r\in\mathcal{M}$ with $X=\bigsqcup_rY_r$ and $\lambda(Y_r),\mu(Y_r)<+\infty$.\\
Apply the finite case to $\lambda|_{Y_r},\mu|_{Y_r}$, get $f_r$, and set $f=\sum_r f_r\mathbf{1}_{Y_r}$. Then $\lambda(E)=\sum_r\lambda(E\cap Y_r)=\int_Ef\,d\mu$.\\
Since $\int_{Y_r}f\,d\mu=\lambda(Y_r)<+\infty$, $f<+\infty$ $\mu$-a.e. If $f,g$ both represent $\lambda$, then $f=g$ $\mu$-a.e. on each $Y_r$ by Proposition \ref{prop:l1-level-sets-null-and-indefinite-integrals}; hence $f=g$ $\mu$-a.e. on $X$.\\
\hlt{Structural Commentary}: Hahn decompositions of $\lambda-c\mu$ build lower step approximations to the desired density.\\
\hlt{Key Insight}: The density is recovered by measuring where $\lambda$ lies between consecutive multiples of $\mu$.
\end{proof}
\end{theorem}

\begin{example} \label{ex:radon-nikodym-needs-sigma-finite-source} \hlt{A Non-$\sigma$-Finite Failure}\\
Let $X=[0,1]$, $\mathcal{M}=\mathcal{B}([0,1])$, $\mu=\#$, and $\lambda=m$. Since $\mu(E)=0\Leftrightarrow E=\emptyset$, $\lambda\ll\mu$.\\
If $\lambda(E)=\int_E f\,d\mu$ for some $f\geq0$, then $0=\lambda(\{x\})=\int_{\{x\}}f\,d\mu=f(x)$ for every $x$, so $\lambda([0,1])=0$, contradiction.\\
Thus RN can fail when $\mu$ is not $\sigma$-finite.
\end{example}

\begin{example} \label{ex:countable-cocountable-radon-nikodym-failure} \hlt{Another Non-$\sigma$-Finite Failure}\\
Let $X$ be uncountable and $\mathcal{M}\defeq\{E\subseteq X:E\text{ is countable or }X\setminus E\text{ is countable}\}$.\\
Set $\mu(E)=\#E$ for finite $E$ and $\mu(E)=+\infty$ otherwise; set $\lambda(E)=0$ for countable $E$ and $\lambda(E)=+\infty$ otherwise. Since $\mu(E)=0\Leftrightarrow E=\emptyset$, $\lambda\ll\mu$. If $\lambda(E)=\int_Ef\,d\mu$ with $f\geq0$, then $0=\lambda(\{x\})=\int_{\{x\}}f\,d\mu=f(x)$ for all $x$, so $\lambda(X)=0$, contradiction.\\
\hlt{Structural Commentary}: This is the countable--cocountable version of Example \ref{ex:radon-nikodym-needs-sigma-finite-source}.\\
\hlt{Key Insight}: Non-$\sigma$-finite references can erase all singleton tests.
\end{example}

\begin{remark} \label{rem:rn-derivative-integrability-finiteness} \hlt{When the Radon--Nikodym Derivative Is Integrable}\\
In Theorem \ref{thm:radon-nikodym-sigma-finite}, $\frac{d\lambda}{d\mu}\in L^1(\mu)\Leftrightarrow\lambda(X)<+\infty$, since $\int_X(d\lambda/d\mu)\,d\mu=\lambda(X)$.
\end{remark}

\begin{theorem} \label{thm:radon-nikodym-extended-dominated-measure} \hlt{Extended-Valued Radon--Nikodym Theorem}\\
Let $\mu$ be $\sigma$-finite, let $\lambda$ be a measure on $(X,\mathcal{M})$, and assume $\lambda\ll\mu$.\\
Then $\exists$ measurable $f:X\to[0,+\infty]$ such that $\lambda(E)=\int_Ef\,d\mu$ for all $E\in\mathcal{M}$; any two such $f$ agree $\mu$-a.e.
\begin{proof}
Assume first $\mu(X)<+\infty$. Let $(P_n,N_n)$ be Hahn for $\lambda-n\mu$, and put $P=\bigcap_{n\geq1}P_n$, $N=\bigcup_{n\geq1}N_n$.\\
If $E\subseteq N_n$, then $\lambda(E)\leq n\mu(E)$; hence $\lambda|_N$ is $\sigma$-finite and Theorem \ref{thm:radon-nikodym-sigma-finite} gives $g:N\to[0,+\infty)$ with $\lambda(E\cap N)=\int_{E\cap N}g\,d\mu$. If $E\subseteq P$, then $\lambda(E)\geq n\mu(E)$ for all $n$, so $\mu(E)>0\Rightarrow\lambda(E)=+\infty$, while $\mu(E)=0\Rightarrow\lambda(E)=0$. Thus $f=g\mathbf{1}_N+\infty\cdot\mathbf{1}_P$ represents $\lambda$.\\
For general $\sigma$-finite $\mu$, choose $X=\bigsqcup_rY_r$ with $\mu(Y_r)<+\infty$, apply the finite-$\mu$ case to each restriction, and set $f=\sum_rf_r\mathbf{1}_{Y_r}$. Then $\lambda(E)=\sum_r\lambda(E\cap Y_r)=\int_Ef\,d\mu$. If $f,h$ represent $\lambda$, then on each $Y_r$, $\int_Af\,d\mu=\int_Ah\,d\mu$ for all $A\subseteq Y_r$; applying this to $\{f<h\}$ and $\{h<f\}$ gives $f=h$ $\mu$-a.e.\\
\hlt{Structural Commentary}: If $\lambda$ is not $\sigma$-finite, the density may have a genuinely infinite part.\\
\hlt{Key Insight}: Hahn decompositions isolate finite and infinite density regions.
\end{proof}
\end{theorem}

\begin{theorem} \label{thm:lebesgue-decomposition-measures} \hlt{Lebesgue Decomposition Theorem}\\
Let $\lambda,\mu$ be $\sigma$-finite measures on $(X,\mathcal{M})$. Then $\exists!$ measures $\lambda_s,\lambda_a$ with $\lambda=\lambda_s+\lambda_a$, $\lambda_s\perp\mu$, and $\lambda_a\ll\mu$.
\begin{proof}
Set $\rho=\lambda+\mu$. Choose $f=d\lambda/d\rho$ and $g=d\mu/d\rho$.\\
Let $A=\{g=0\}$, $B=\{g>0\}$, $\lambda_s(E)\defeq\lambda(E\cap A)$, and $\lambda_a(E)\defeq\lambda(E\cap B)$. Since $\mu(A)=\int_Ag\,d\rho=0$, $\lambda_s\perp\mu$.\\
If $\mu(E)=0$, then $g=0$ $\rho$-a.e. on $E$, so $\rho(E\cap B)=0$ and $\lambda_a(E)=0$; hence $\lambda_a\ll\mu$.\\
For uniqueness, let $(\tilde{\lambda}_s,\tilde{\lambda}_a)$ be another decomposition.\\
Since $\lambda_s\perp\mu$ and $\tilde{\lambda}_a\ll\mu$, also $\lambda_s\perp\tilde{\lambda}_a$; choose $X=C\cup D$ with $\lambda_s(C)=0$ and $\tilde{\lambda}_a(D)=0$.\\
Then $\lambda_s(E)=\lambda_s(E\cap D)\leq\lambda(E\cap D)=\tilde{\lambda}_s(E\cap D)\leq\tilde{\lambda}_s(E)$. By symmetry, $\lambda_s=\tilde{\lambda}_s$, hence $\lambda_a=\tilde{\lambda}_a$.\\
\hlt{Structural Commentary}: Every measure decomposes into a $\mu$-visible part and a part living on a $\mu$-null set.\\
\hlt{Key Insight}: Differentiate both measures with respect to their sum, then split where the $\mu$-density vanishes.
\end{proof}
\end{theorem}

\begin{figure}[H]
  \centering
  \begin{tikzpicture}[
  scale=0.88,
  every node/.style={font=\scriptsize},
  box/.style={draw=black!50, line width=0.6pt, align=center, minimum width=1.62cm, minimum height=0.58cm, fill=white},
  formula/.style={box, minimum width=2.42cm, minimum height=0.70cm},
  posbox/.style={box, fill=blue!18},
  neg/.style={box, fill=red!10},
  ac/.style={box, fill=teal!12},
  sing/.style={box, fill=orange!16},
  arrow/.style={-{Latex[length=2mm]}, line width=0.55pt, black!60}
]
  \node[box] (charge) at (0,1.25) {charge\\$\lambda$};
  \node[posbox] (plus) at (-1.32,0.18) {$\lambda^+$\\positive};
  \node[neg] (minus) at (1.32,0.18) {$\lambda^-$\\negative};
  \node[formula] (jordan) at (0,-0.98) {$\lambda=\lambda^+-\lambda^-$\\$\abs{\lambda}=\lambda^++\lambda^-$};

  \draw[arrow] (charge) -- (plus);
  \draw[arrow] (charge) -- (minus);
  \draw[arrow] (plus) -- (jordan);
  \draw[arrow] (minus) -- (jordan);

  \node[box] (measure) at (4.60,1.25) {measure\\$\lambda$};
  \node[sing] (singular) at (3.28,0.18) {$\lambda_s\perp\mu$\\singular};
  \node[ac] (absolute) at (5.92,0.18) {$\lambda_a\ll\mu$\\visible};
  \node[formula] (leb) at (4.60,-0.98) {$\lambda=\lambda_s+\lambda_a$\\$\lambda_a(E)=\int_E f\,d\mu$};

  \draw[arrow] (measure) -- (singular);
  \draw[arrow] (measure) -- (absolute);
  \draw[arrow] (singular) -- (leb);
  \draw[arrow] (absolute) -- (leb);

  \draw[black!25, line width=0.5pt] (2.30,-1.48) -- (2.30,1.72);
\end{tikzpicture}

  \caption{Charges split by sign; measures split into $\lambda_s\perp\mu$ and $\lambda_a\ll\mu$ parts.}
  \label{fig:measure-decomposition-schematic}
\end{figure}

\begin{proposition} \label{prop:radon-nikodym-change-of-measure-calculus} \hlt{Radon--Nikodym Calculus}\\
Let all measures below be $\sigma$-finite.
\begin{enumerate}[label=(\roman*)]
\setlength{\itemsep}{0pt}
\item $\lambda\ll\mu$, $g\geq0$ measurable $\Rightarrow\int_Xg\,d\lambda=\int_Xg(d\lambda/d\mu)\,d\mu$.
\item $\nu\ll\lambda\ll\mu\Rightarrow d\nu/d\mu=(d\nu/d\lambda)(d\lambda/d\mu)$ $\mu$-a.e.
\item $\lambda_1,\lambda_2\ll\mu\Rightarrow d(\lambda_1+\lambda_2)/d\mu=d\lambda_1/d\mu+d\lambda_2/d\mu$ $\mu$-a.e.
\item $\lambda\ll\mu\ll\lambda\Rightarrow d\lambda/d\mu=1/(d\mu/d\lambda)$ $\mu$-a.e.
\end{enumerate}
\begin{proof}
\begin{enumerate}[label=(\roman*)]
\setlength{\itemsep}{0pt}
\item For $g=\mathbf{1}_E$, $\int_Xg\,d\lambda=\lambda(E)=\int_E(d\lambda/d\mu)\,d\mu=\int_Xg(d\lambda/d\mu)\,d\mu$.\\
By linearity the identity holds for nonnegative simple $g$; for measurable $g\geq0$, take simple $g_n\uparrow g$ and apply monotone convergence to both sides.
\item Let $h\defeq(d\nu/d\lambda)(d\lambda/d\mu)$.\\
For every $E\in\mathcal{M}$, (i) with $g=\mathbf{1}_E(d\nu/d\lambda)$ gives $\nu(E)=\int_E(d\nu/d\lambda)\,d\lambda=\int_Eh\,d\mu$.\\
Hence $h$ represents $\nu$ relative to $\mu$, so $h=d\nu/d\mu$ $\mu$-a.e.
\item Let $h\defeq d\lambda_1/d\mu+d\lambda_2/d\mu$. For every $E\in\mathcal{M}$, $\int_Eh\,d\mu=\lambda_1(E)+\lambda_2(E)=(\lambda_1+\lambda_2)(E)$.\\
Hence $h$ represents $\lambda_1+\lambda_2$ relative to $\mu$, so uniqueness gives the identity.
\item Apply (ii) to $\mu\ll\lambda\ll\mu$. Since $d\mu/d\mu=1$, $1=(d\mu/d\lambda)(d\lambda/d\mu)$ $\mu$-a.e.\\
Also $\mu\ll\lambda$ gives $d\mu/d\lambda>0$ $\mu$-a.e.; divide to obtain $d\lambda/d\mu=1/(d\mu/d\lambda)$ $\mu$-a.e.
\end{enumerate}
\hlt{Structural Commentary}: Radon--Nikodym derivatives obey chain and linearity rules.\\
\hlt{Key Insight}: Prove identities by integrating over arbitrary measurable sets, then invoke uniqueness.
\end{proof}
\end{proposition}

\begin{corollary} \label{cor:absolute-continuous-singular-zero-measure} \hlt{Absolutely Continuous and Singular Forces Zero}\\
For measures $\lambda,\mu$, $\lambda\ll\mu$ and $\lambda\perp\mu$ imply $\lambda=0$. For a charge $\eta$, $\abs{\eta}\perp\mu$ implies $\eta^+\perp\mu$ and $\eta^-\perp\mu$.
\begin{proof}
If $\lambda\perp\mu$, choose $X=A\cup B$ with $\lambda(A)=0$ and $\mu(B)=0$. Since $\lambda\ll\mu$, $\lambda(B)=0$, hence $\lambda(X)=0$.\\
If $\abs{\eta}\perp\mu$, choose $X=A\cup B$ with $\abs{\eta}(A)=0$ and $\mu(B)=0$. Since $\eta^\pm\leq\abs{\eta}$, $\eta^\pm(A)=0$, so $\eta^\pm\perp\mu$.\\
\hlt{Structural Commentary}: Absolute continuity and singular support are incompatible except at zero.\\
\hlt{Key Insight}: Visible and singular forces zero.
\end{proof}
\end{corollary}

\begin{corollary} \label{cor:radon-nikodym-for-charges} \hlt{Radon--Nikodym for Charges}\\
Let $\mu$ be $\sigma$-finite and let $\lambda$ be a charge with $\lambda\ll\mu$. Then $\exists f\in\mathscr{L}^1(\mu)$ such that, for all $E\in\mathcal{M}$, $\lambda(E)=\int_E f\,d\mu$, $\lambda^+(E)=\int_E f^+\,d\mu$, $\lambda^-(E)=\int_E f^-\,d\mu$, and $\abs{\lambda}(E)=\int_E\abs{f}\,d\mu$.
\begin{proof}
Since $\abs{\lambda}\ll\mu$, $\lambda^\pm\ll\mu$. By Radon--Nikodym, choose $f_\pm\geq0$ with $\lambda^\pm(E)=\int_Ef_\pm\,d\mu$. Since $\lambda^+\perp\lambda^-$, $f_+f_-=0$ $\mu$-a.e. Set $f=f_+-f_-$. Then $\int_X\abs{f}\,d\mu=\int_X(f_++f_-)\,d\mu=\abs{\lambda}(X)<+\infty$, so $f\in\mathscr{L}^1(\mu)$, $f^+=f_+$, $f^-=f_-$ $\mu$-a.e., and $\lambda(E)=\lambda^+(E)-\lambda^-(E)=\int_Ef\,d\mu$.\\
\hlt{Structural Commentary}: Signed absolute continuity is reduced to the two positive Jordan pieces.\\
\hlt{Key Insight}: Differentiate $\lambda^+$ and $\lambda^-$ separately, then subtract.
\end{proof}
\end{corollary}

\begin{example} \label{ex:stieltjes-basic-generators-rn-classification} \hlt{Radon--Nikodym Classification of Basic Stieltjes Generators}\\
Let $m$ be Lebesgue measure on $\mathcal{B}(\R)$, and let $\lambda_i$ be the Borel--Stieltjes measure generated by $g_i$.
\begin{enumerate}[label=(\roman*)]
\setlength{\itemsep}{0pt}
\item $g_1(x)=2x$: $\lambda_1=2m$, so $\lambda_1\ll m$, $d\lambda_1/dm=2$, $m\ll\lambda_1$, $dm/d\lambda_1=1/2$, and $\lambda_1(\R)=+\infty$.
\item $g_2(x)=\arctan x$: $\lambda_2(E)=\int_E(1+x^2)^{-1}\,dm$, so $d\lambda_2/dm=(1+x^2)^{-1}$, $\lambda_2(\R)=\pi$, $m\ll\lambda_2$, and $dm/d\lambda_2=1+x^2$.
\item $g_3=\mathbf{1}_{[0,+\infty)}$: $\lambda_3=\delta_0$,\\
so $\lambda_3\perp m$, $\lambda_3(\R)=1$, $\lambda_3\not\ll m$, and $m\not\ll\lambda_3$.
\item $g_4(x)=x\mathbf{1}_{[0,+\infty)}(x)$: $\lambda_4(E)=m(E\cap[0,+\infty))$, so $\lambda_4\ll m$, $d\lambda_4/dm=\mathbf{1}_{[0,+\infty)}$, $m\not\ll\lambda_4$, and $\lambda_4(\R)=+\infty$.
\end{enumerate}
Thus $\lambda_1,\lambda_2,\lambda_4\ll m$, $\lambda_3\perp m$, only $\lambda_2,\lambda_3$ are finite, and $m\ll\lambda_i$ exactly for $i=1,2$.\\
\hlt{Structural Commentary}: Stieltjes generators split naturally into density, atom, and restriction examples.\\
\hlt{Key Insight}: A positive density makes null sets equivalent; an atom lives on a Lebesgue-null set.
\end{example}

\begin{theorem} \label{thm:charges-total-variation-banach} \hlt{Charges Form a Banach Space}\\
Let $\operatorname{ca}(X,\mathcal{M})\defeq\{\lambda:\mathcal{M}\to\R:\lambda\text{ is a charge}\}$ and $\norm{\lambda}_{\mathrm{var}}\defeq\abs{\lambda}(X)$.\\
With pointwise operations, $(\operatorname{ca}(X,\mathcal{M}),\norm{\cdot}_{\mathrm{var}})$ is Banach.
\begin{proof}
Linearity follows from charge linearity.\\
Definiteness follows from $\abs{\lambda}(X)=0\Rightarrow\abs{\lambda}(E)=0\ \forall E$, homogeneity from $\abs{a\lambda}=|a|\abs{\lambda}$, and the triangle inequality from $\sum_j\abs{(\lambda+\eta)(A_j)}\leq\sum_j\abs{\lambda(A_j)}+\sum_j\abs{\eta(A_j)}$ over finite measurable partitions $(A_j)$.\\
Let $(\lambda_n)$ be variation-Cauchy. Since $\abs{\lambda_n(E)-\lambda_m(E)}\leq\norm{\lambda_n-\lambda_m}_{\mathrm{var}}$, define $\lambda(E)\defeq\lim_n\lambda_n(E)$; then $\lambda$ is finitely additive. For fixed $n$ and finite disjoint $(A_j)$, $\sum_j\abs{(\lambda-\lambda_n)(A_j)}\leq\limsup_m\norm{\lambda_m-\lambda_n}_{\mathrm{var}}$.\\
For disjoint $(E_k)$, set $E=\bigcup_kE_k$ and $R_K=\bigcup_{k>K}E_k$. Choose $n$ with the preceding limsup $<\varepsilon$.\\
Since $\lambda_n(R_K)\to0$, for large $K$, $\abs{\lambda(R_K)}\leq\abs{\lambda_n(R_K)}+\abs{(\lambda-\lambda_n)(R_K)}<2\varepsilon$. Hence $\lambda(R_K)\to0$, so $\lambda$ is countably additive. Taking suprema over finite partitions gives $\norm{\lambda-\lambda_n}_{\mathrm{var}}\leq\limsup_m\norm{\lambda_m-\lambda_n}_{\mathrm{var}}\to0$.\\
\hlt{Structural Commentary}: Total variation is the signed-measure analogue of an $L^1$ norm.\\
\hlt{Key Insight}: Variation-Cauchy controls every finite partition.
\end{proof}
\end{theorem}
